\documentclass[11pt]{article}
\usepackage[margin=1in]{geometry}
\usepackage{amsmath,amssymb,amsthm}
\usepackage{hyperref}
\usepackage{booktabs}
\usepackage{listings}
\usepackage{xcolor}

\lstset{
  language=Python,
  basicstyle=\ttfamily\small,
  keywordstyle=\color{blue},
  commentstyle=\color{gray},
  stringstyle=\color{red!70!black},
  showstringspaces=false,
  frame=single,
  breaklines=true,
}

\newtheorem{theorem}{Theorem}
\newtheorem{lemma}[theorem]{Lemma}

\title{Analysis of the Suzuki Recursion Formula\\in the QDK Trotter Implementation}
\author{QDK-Chemistry}
\date{\today}

\begin{document}
\maketitle

\section{Background: Suzuki's Fractal Decomposition}

Consider a Hamiltonian $H = \sum_{j=1}^{L} H_j$.  The goal of Trotter--Suzuki
product formulas is to approximate the time-evolution operator $e^{-iHt}$ by a
product of exponentials of the individual terms $H_j$.

\subsection{Second-order formula (Strang splitting)}

The second-order symmetric (Strang) splitting is:
\begin{equation}\label{eq:S2}
  S_2(t) = e^{-iH_1 t/2}\, e^{-iH_2 t/2}\, \cdots\, e^{-iH_{L-1} t/2}\,
            e^{-iH_L t}\,
            e^{-iH_{L-1} t/2}\, \cdots\, e^{-iH_1 t/2}.
\end{equation}
This satisfies
\begin{equation}
  \bigl\lVert e^{-iHt} - S_2(t) \bigr\rVert = \mathcal{O}(t^3).
\end{equation}

\subsection{Higher-order formulas via recursion}

Suzuki~\cite{Suzuki1990,Suzuki1991} showed that higher-order approximations can
be constructed recursively.  Given an approximation $S_{2k}(t)$ of order $2k$
(i.e., error $\mathcal{O}(t^{2k+1})$), one obtains an approximation of order
$2k+2$ via:
\begin{equation}\label{eq:recursion}
  S_{2k+2}(t)
  = \bigl[S_{2k}(p_k\, t)\bigr]^2\;
    S_{2k}\!\bigl((1 - 4p_k)\, t\bigr)\;
    \bigl[S_{2k}(p_k\, t)\bigr]^2,
\end{equation}
where the parameter $p_k$ is chosen to cancel the leading-order error:
\begin{equation}\label{eq:pk}
  \boxed{p_k = \frac{1}{4 - 4^{1/(2k+1)}}}.
\end{equation}

\begin{theorem}[Suzuki 1991]\label{thm:suzuki}
If $S_{2k}(t) = e^{-iHt} + \mathcal{O}(t^{2k+1})$, then
$S_{2k+2}(t)$ defined by~\eqref{eq:recursion} with $p_k$ from~\eqref{eq:pk}
satisfies $S_{2k+2}(t) = e^{-iHt} + \mathcal{O}(t^{2k+3})$.
\end{theorem}

\section{Constructing the Fourth-Order Formula}

To build the fourth-order formula $S_4$ from the second-order formula $S_2$, we
apply one level of Suzuki recursion with $k = 1$ (promoting order $2 \cdot 1 = 2$
to order $2 \cdot 2 = 4$):
\begin{equation}\label{eq:S4}
  S_4(t) = \bigl[S_2(p_1\, t)\bigr]^2\;
           S_2\!\bigl((1-4p_1)\, t\bigr)\;
           \bigl[S_2(p_1\, t)\bigr]^2,
\end{equation}
with
\begin{equation}\label{eq:p1}
  p_1 = \frac{1}{4 - 4^{1/(2 \cdot 1 + 1)}} = \frac{1}{4 - 4^{1/3}}.
\end{equation}

\noindent
Numerically, $p_1 = (4 - 4^{1/3})^{-1} \approx 0.41449$.

\section{The Bug in the QDK Implementation}

The QDK's \texttt{suzuki\_recursion} function in
\texttt{source/pip/qsharp/applications/magnets/trotter/trotter.py}
implements the recursion as follows:

\begin{lstlisting}
def suzuki_recursion(trotter: TrotterStep) -> TrotterStep:
    # ...
    p = 1 / (4 - 4 ** (1 / (2 * trotter.order + 1)))
    # ...
\end{lstlisting}

When called on a Strang splitting (\texttt{trotter.order = 2}), this computes:
\begin{equation}\label{eq:qdk_p}
  p_{\text{QDK}} = \frac{1}{4 - 4^{1/(2 \cdot 2 + 1)}} = \frac{1}{4 - 4^{1/5}}
  \approx 0.37307.
\end{equation}

\subsection{What went wrong}

The \texttt{trotter.order} attribute equals $2$ (the order of the input
formula).  In Suzuki's notation, the input is $S_{2k}$ with $2k = 2$, so $k =
1$.  The correct parameter is:
\begin{equation}
  p_k = p_1 = \frac{1}{4 - 4^{1/(2 \cdot 1 + 1)}} = \frac{1}{4 - 4^{1/3}}.
\end{equation}

The code substitutes \texttt{trotter.order} (which is $2k = 2$) directly in
place of $k$:
\begin{equation}
  p_{\text{QDK}} = \frac{1}{4 - 4^{1/(2 \cdot (2k) + 1)}} = \frac{1}{4 - 4^{1/(4k+1)}}.
\end{equation}

For $k=1$ this gives $p_{\text{QDK}} = 1/(4-4^{1/5})$ instead of the correct
$p_1 = 1/(4-4^{1/3})$.

\subsection{Correct fix}

The fix is to use $k = \texttt{trotter.order} / 2$:
\begin{lstlisting}
k = trotter.order // 2   # order=2 -> k=1
p = 1 / (4 - 4 ** (1 / (2 * k + 1)))
\end{lstlisting}

\noindent
Or equivalently, use \texttt{trotter.order - 1} as the exponent denominator:
\begin{lstlisting}
p = 1 / (4 - 4 ** (1 / (trotter.order + 1)))
\end{lstlisting}

\noindent
since $2k + 1 = \texttt{trotter.order} + 1$ when $\texttt{trotter.order} = 2k$.

\section{Impact: Loss of Convergence Order}

With the wrong $p$, the leading-order error term in $S_{2k}$ is \emph{not}
cancelled by the recursion.  The resulting formula still converges, but only at
the original order $2k$ rather than the intended $2k+2$.

\subsection{Numerical verification}

We verify this with a 2-qubit Ising Hamiltonian
$H = Z_0 Z_1 + 0.5\,(X_0 + X_1)$, computing $e^{-iHt}$ with $t = 1$ using
$N$ Trotter steps.  The error is
$\varepsilon(N) = \lVert U_{\text{Trotter}} - U_{\text{exact}} \rVert$.

\begin{table}[h]
\centering
\begin{tabular}{rcccc}
\toprule
$N$ & $\varepsilon_{\text{standard}}$ & order & $\varepsilon_{\text{QDK}}$ & order \\
\midrule
 4  & $5.95 \times 10^{-5}$ & ---  & $2.12 \times 10^{-3}$ & ---  \\
 8  & $3.70 \times 10^{-6}$ & 4.01 & $5.34 \times 10^{-4}$ & 1.99 \\
16  & $2.31 \times 10^{-7}$ & 4.00 & $1.34 \times 10^{-4}$ & 2.00 \\
32  & $1.44 \times 10^{-8}$ & 4.00 & $3.34 \times 10^{-5}$ & 2.00 \\
64  & $9.02 \times 10^{-10}$& 4.00 & $8.35 \times 10^{-6}$ & 2.00 \\
\bottomrule
\end{tabular}
\caption{Convergence order $= \log_2(\varepsilon_N / \varepsilon_{2N})$.  The
standard formula achieves 4th-order convergence; the QDK formula achieves only
2nd-order.}
\label{tab:convergence}
\end{table}

\noindent
Table~\ref{tab:convergence} confirms that the standard formula converges as
$\mathcal{O}(N^{-4})$ (4th~order), while the QDK formula converges as
$\mathcal{O}(N^{-2})$ (2nd~order).  The Suzuki recursion with the wrong $p$
does not achieve the intended order promotion.

\section{Verification of the qdk-chemistry Implementation}

We verify that the qdk-chemistry Trotter implementation achieves the correct
convergence order by running the full pipeline end-to-end:
\texttt{create\_ising\_hamiltonian} $\to$ \texttt{Trotter.run()} $\to$
\texttt{to\_circuit()} $\to$ Cirq unitary extraction $\to$ comparison with
$e^{-iHt}$.

\subsection{2-qubit Ising chain}

$H = Z_0 Z_1 + 0.5\,X_0 + 0.5\,X_1$, $t = 1$.

\begin{table}[h]
\centering
\begin{tabular}{rcccc}
\toprule
$N$ & $\varepsilon$ (order~2) & conv.\ order & $\varepsilon$ (order~4) & conv.\ order \\
\midrule
  2 & $1.01 \times 10^{-1}$ & --- & $9.74 \times 10^{-4}$ & --- \\
  4 & $2.44 \times 10^{-2}$ & 2.04 & $5.95 \times 10^{-5}$ & 4.03 \\
  8 & $6.06 \times 10^{-3}$ & 2.01 & $3.70 \times 10^{-6}$ & 4.01 \\
 16 & $1.51 \times 10^{-3}$ & 2.00 & $2.31 \times 10^{-7}$ & 4.00 \\
 32 & $3.78 \times 10^{-4}$ & 2.00 & $1.44 \times 10^{-8}$ & 4.00 \\
 64 & $9.45 \times 10^{-5}$ & 2.00 & $9.02 \times 10^{-10}$ & 4.00 \\
128 & $2.36 \times 10^{-5}$ & 2.00 & $5.64 \times 10^{-11}$ & 4.00 \\
\bottomrule
\end{tabular}
\caption{qdk-chemistry convergence for the 2-qubit Ising chain.
Both 2nd- and 4th-order formulas achieve their expected convergence rates.}
\end{table}

\subsection{4-qubit Ising lattice ($2\times 2$)}

$H = \sum_{\langle i,j\rangle} Z_i Z_j + 0.5 \sum_i X_i$ on a $2 \times 2$
square lattice (8 Pauli terms), $t = 1$.

\begin{table}[h]
\centering
\begin{tabular}{rcccc}
\toprule
$N$ & $\varepsilon$ (order~2) & conv.\ order & $\varepsilon$ (order~4) & conv.\ order \\
\midrule
 2 & $4.38 \times 10^{-1}$ & --- & $9.23 \times 10^{-3}$ & --- \\
 4 & $1.03 \times 10^{-1}$ & 2.08 & $5.65 \times 10^{-4}$ & 4.03 \\
 8 & $2.55 \times 10^{-2}$ & 2.02 & $3.54 \times 10^{-5}$ & 4.00 \\
16 & $6.35 \times 10^{-3}$ & 2.00 & $2.22 \times 10^{-6}$ & 4.00 \\
32 & $1.59 \times 10^{-3}$ & 2.00 & $1.39 \times 10^{-7}$ & 4.00 \\
64 & $3.97 \times 10^{-4}$ & 2.00 & $8.67 \times 10^{-9}$ & 4.00 \\
\bottomrule
\end{tabular}
\caption{qdk-chemistry convergence for the $2\times 2$ Ising lattice
(with \texttt{optimize\_term\_ordering=True}).}
\end{table}

\noindent
All tests confirm that the qdk-chemistry implementation achieves the correct
convergence order for both 2nd- and 4th-order Trotter formulas.

\section{Summary}

\begin{itemize}
  \item The Suzuki recursion \eqref{eq:recursion} promotes order $2k$ to
    $2k+2$ using $p_k = 1/(4 - 4^{1/(2k+1)})$, where $k$ indexes the
    \emph{input} formula order as $2k$.
  \item The QDK code uses \texttt{trotter.order} (which equals $2k$, not $k$)
    directly in the formula, computing
    $p = 1/(4 - 4^{1/(2 \cdot 2k + 1)}) = 1/(4 - 4^{1/(4k+1)})$.
  \item For the $S_2 \to S_4$ promotion ($k=1$), this gives
    $1/(4-4^{1/5}) \approx 0.373$ instead of the correct
    $1/(4-4^{1/3}) \approx 0.414$.
  \item Numerical experiments confirm the resulting formula is only 2nd-order
    accurate, not 4th-order as intended.
\end{itemize}

\begin{thebibliography}{9}
\bibitem{Suzuki1990}
  M.~Suzuki,
  ``Fractal decomposition of exponential operators with applications to
  many-body theories and Monte Carlo simulations,''
  \textit{Phys.\ Lett.\ A}, vol.~146, pp.~319--323, 1990.

\bibitem{Suzuki1991}
  M.~Suzuki,
  ``General theory of fractal path integrals with applications to many-body
  theories and statistical physics,''
  \textit{J.\ Math.\ Phys.}, vol.~32, pp.~400--407, 1991.
\end{thebibliography}

\end{document}
